\documentclass[11pt,a4paper]{article}

\usepackage[margin=1in]{geometry}
\usepackage{amsmath,amssymb,amsthm,mathtools}
\usepackage{booktabs}
\usepackage{enumitem}
\usepackage[hidelinks]{hyperref}

\newtheorem{theorem}{Theorem}[section]
\newtheorem{proposition}[theorem]{Proposition}
\newtheorem{lemma}[theorem]{Lemma}
\newtheorem{corollary}[theorem]{Corollary}
\theoremstyle{definition}
\newtheorem{definition}[theorem]{Definition}
\theoremstyle{remark}
\newtheorem{remark}[theorem]{Remark}

\newcommand{\R}{\mathbb{R}}
\newcommand{\Om}{\widetilde{\Omega}}
\newcommand{\Sph}{\mathbb{S}^{N-1}}
\newcommand{\norm}[1]{\lVert#1\rVert}
\newcommand{\abs}[1]{\lvert#1\rvert}
\newcommand{\eps}{\varepsilon}
\newcommand{\dist}{\operatorname{dist}}
\newcommand{\Lip}{\operatorname{Lip}}
\newcommand{\bdry}{\partial}
\newcommand{\sph}{\mathrm{sph}}

\title{Schauder bootstrap and Hölder interpolation:\\
       a quantitative $C^{1,\gamma}$ $\to$ small $C^{2,\gamma_{0}}$ module\\
       for sharp Faber--Krahn stability}
\author{Plan~1 / Route A, building block \textsc{SchauderInterpolation}}
\date{}

\begin{document}
\maketitle

\begin{abstract}
This note isolates the quantitative regularity step that turns the global
$C^{1,\gamma}$ graph produced by \texttt{GraphEntry} into a graph of
\emph{small} $C^{2,\gamma_{0}}$ norm --- the smallness hypothesis required
by BDV's nearly spherical theorem (Theorem~3.3 of~\cite{BDV}). The argument
is constructive: no compactness, no contradiction. All constants are named
with their $(N,R)$-dependence. The two black boxes are
Kinderlehrer--Nirenberg~\cite{KN}~\S2--\S3 (hodograph straightening of the
torsion problem) and Lieberman~\cite{Li86} (Schauder estimates for oblique
problems with Hölder coefficients).
\end{abstract}

\tableofcontents

\section{Input: what \texttt{GraphEntry} and BDV supply}\label{sec:input}

We work under the standing setup of Route~A. After single-set selection
and volume normalization
($r_{*}=(|U_{*}|/|B_{1}|)^{1/N}$, $|r_{*}-1|\le C\delta_{T}$), the
rescaled minimizer $\Om\subset B_{R}$ has $|\Om|=|B_{1}|$.

\subsection*{(a)--(b) Graph and Hausdorff size.}
\texttt{GraphEntry} produces $g:\Sph\to\R$ with
\begin{equation}\label{eq:graph}
   \bdry\Om=\bigl\{(1+g(\theta))\theta:\theta\in\Sph\bigr\},
\end{equation}
satisfying
\begin{align}
   \norm{g}_{C^{1,\gamma}(\Sph)}    & \le M_{1}(N,R),
   \label{eq:input-a}\\
   \norm{g}_{L^{\infty}(\Sph)}      & \le C_{H}''(N,R)\,
                                      \alpha(\Om)^{1/(2N)},
   \label{eq:input-b}
\end{align}
where $\alpha(\Om)$ is the BDV smooth asymmetry, $\gamma=\gamma(N,R)\in(0,1)$
is the Alt--Caffarelli exponent of BDV Theorem~4.18, $M_{1}(N,R)$ comes
from the flatness theorem together with BDV Lemma~4.16, and $C_{H}''(N,R)$
is the asymmetry-to-Hausdorff constant of \texttt{GraphEntry}.

\subsection*{(c)--(d) BDV regularity package.}
The torsion potential $\widetilde u\in W^{1,\infty}_{0}(\Om)\cap C^{\infty}(\Om)$
solves $-\Delta\widetilde u=1$ in $\Om$, $\widetilde u=0$ on $\bdry\Om$. We
use:
\begin{enumerate}[label=\textup{(\roman*)},leftmargin=2.5em]
  \item \emph{Lipschitz nondegeneracy} (BDV Lemma~4.12):
        $L_{u}^{-1}\dist(x,\bdry\Om)\le\widetilde u(x)\le L_{u}\dist(x,\bdry\Om)$
        for $x\in\Om\cap\{\dist(\cdot,\bdry\Om)\le r_{d}\}$,
        with $L_{u}=L_{u}(N,R)$, $r_{d}=r_{d}(N,R)$.
  \item \emph{Smooth Bernoulli coefficient} (BDV Lemma~4.16): the
        coefficient $q(y):=|\nabla\widetilde u(y)|$ on $\bdry\Om$ admits a
        $C^{1,\gamma}$ extension to a one-sided collar with
        $q_{-}\le q\le q_{+}$ and $\norm{q}_{C^{1,\gamma}}\le L_{q}(N,R)$.
\end{enumerate}
All constants depend only on $(N,R)$ and degrade by at most $1+O(\delta_{T})$
under volume normalization.

\section{Hodograph straightening following Kinderlehrer--Nirenberg}\label{sec:hodo}

This section reduces the curved obstacle problem on $\Om$ to a
fixed-domain oblique BVP, following \cite{KN}~\S2--\S3.

\subsection{Chart, partial hodograph, inverse map}\label{ss:hodo-chart}

Fix $y_{0}\in\bdry\Om$. By~\eqref{eq:input-a} the inward normal $\nu(y_{0})$
makes an angle $\le\arctan M_{1}$ with the radial direction. After a
rigid motion assume $y_{0}=0$, $\nu(y_{0})=e_{N}$, and write $x=(x',x_{N})$.
By~\eqref{eq:input-a} there is $R_{0}=R_{0}(N,R)>0$ such that in
$\mathcal{C}_{R_{0}}=\{|x'|<R_{0},|x_{N}|<R_{0}\}$,
\[
   \bdry\Om\cap\mathcal{C}_{R_{0}}
      =\{x_{N}=\widetilde g(x'):|x'|<R_{0}\},
   \qquad \norm{\widetilde g}_{C^{1,\gamma}}\le M_{1}'(N,R),
\]
and $\Om\cap\mathcal{C}_{R_{0}}=\{x_{N}>\widetilde g(x')\}$. Set
$u(x):=\widetilde u(x',x_{N})$; then
\begin{equation}\label{eq:torsion-cyl}
   -\Delta u=1\text{ in }\{x_{N}>\widetilde g(x')\},\qquad
   u=0\text{ on }\{x_{N}=\widetilde g(x')\},
\end{equation}
together with the Bernoulli condition $|\nabla u|=q$ on the free boundary,
$q\in C^{1,\gamma}$, $q\ge q_{-}>0$ (BDV Lemma~4.16).

\paragraph{Partial hodograph.}
By Lipschitz nondegeneracy and positivity of $q$, in a one-sided
neighborhood of $0$ we have $\partial_{N}u(x)\ge q_{-}/2>0$. The
partial hodograph map
\[
   \Phi:(x',x_{N})\mapsto(x',\,t),\qquad t:=u(x',x_{N}),
\]
is a $C^{1,\gamma}$ diffeomorphism from $\Om\cap\mathcal{C}_{R_{0}/2}$
onto a domain $\{(x',t):|x'|<R_{0}/2,0<t<T_{0}\}$ for $T_{0}=T_{0}(N,R)>0$.
Its inverse coordinate gives
\[
   \psi(x',t):=x_{N}=\Phi^{-1}(x',t)_{N},\qquad
   \psi(x',0)=\widetilde g(x'),
\]
exactly the construction of \cite{KN}~\S2.

\subsection{Derived quasilinear PDE}\label{ss:hodo-PDE}

A direct change of variables (\cite{KN} (2.10)--(2.12); paralleling
(4.10)--(4.12) for the free-boundary application) gives
\begin{equation}\label{eq:psi-PDE}
   F[\psi]
     :=\frac{1+|\nabla_{x'}\psi|^{2}}{(\partial_{t}\psi)^{2}}\partial_{t}^{2}\psi
        -2\,\frac{\nabla_{x'}\psi\cdot\nabla_{x'}\partial_{t}\psi}{\partial_{t}\psi}
        +\Delta_{x'}\psi
     =\partial_{t}\psi
\end{equation}
in $\{0<t<T_{0}\}$. \eqref{eq:psi-PDE} is uniformly elliptic in $\psi$ as
long as $\partial_{t}\psi\ge\eta>0$ and $|\nabla_{x'}\psi|\le K$, with
ellipticity constants bounded between $\eta_{*}(\eta,K)$ and
$\eta_{*}^{-1}$.

\subsection{Nonlinear oblique boundary condition}\label{ss:hodo-BC}

The free-boundary condition $|\nabla u|^{2}=q^{2}$ at $t=0$ translates
into
\begin{equation}\label{eq:BC-psi}
   \partial_{t}\psi(x',0)
   = \sqrt{\frac{1+|\nabla_{x'}\psi(x',0)|^{2}}{q(x',\psi(x',0))^{2}}}
   \qquad\text{on }\{t=0\},
\end{equation}
the formula in \cite{KN}~(3.4)--(3.5).

\paragraph{Transversality.}
Linearising $B[\psi]:=\partial_{t}\psi-G(x',\nabla_{x'}\psi,\psi)=0$
gives $\partial B/\partial(\partial_{t}\psi)\equiv1$. In Lieberman's
conventions \cite[\S1]{Li86}, the obliqueness coefficient is
$\beta_{N}=1$, uniformly positive. Tangential coefficients
$\beta_{i}=-\partial G/\partial(\partial_{i}\psi)$, $i<N$, are bounded
by $(N,R,L_{q})$-constants.

\subsection{Initial regularity of $\psi$ and norm equivalence}\label{ss:hodo-norm}

The graph $g$ (spherical chart) and the Cartesian $\widetilde g$ are
related by a $(N,R)$-bilipschitz $C^{\infty}$ change of coordinates, so
\begin{equation}\label{eq:norm-equiv-g-tildeg}
   c(N,R,k,\beta)\norm{g}_{C^{k,\beta}(\Sph)}
   \le\norm{\widetilde g}_{C^{k,\beta}}
   \le C(N,R,k,\beta)\norm{g}_{C^{k,\beta}(\Sph)}.
\end{equation}
$\widetilde g(x')=\psi(x',0)$, so
\begin{equation}\label{eq:norm-equiv-tildeg-psi}
   \norm{\widetilde g}_{C^{k,\beta}(|x'|<R_{0}/2)}
   \le\norm{\psi}_{C^{k,\beta}(\{t=0\})}
   \le\norm{\psi}_{C^{k,\beta}(\{0\le t\le T_{0}/2\})}.
\end{equation}
Together,
\begin{equation}\label{eq:norm-three-way}
   \norm{g}_{C^{k,\beta}(\Sph)}\sim_{N,R}
   \norm{\widetilde g}_{C^{k,\beta}}\sim_{N,R}
   \norm{\psi(\cdot,0)}_{C^{k,\beta}}.
\end{equation}
The initial bound $\norm{\psi}_{C^{1,\gamma}}\le C(N,R)$ on
$\{0\le t\le T_{0}/2\}$ follows from~(i), implicit function theorem on
$t=u(x',x_{N})$, and chain rule. A cover of $\bdry\Om$ by
$N_{0}(N,R)$ charts extends~\eqref{eq:norm-three-way} globally.

\section{Schauder bootstrap via Lieberman's oblique theorem}\label{sec:schauder}

\begin{theorem}[Lieberman \cite{Li86}, Theorem~1]\label{thm:lieberman}
Let $v$ solve
$a^{ij}(x,v,\nabla v)\partial_{ij}v+b(x,v,\nabla v)=0$ in $B^{+}_{R_{0}/2}$,
$b^{i}(x,v,\nabla v)\partial_{i}v+b_{0}(x,v)=0$ on $\{x_{N}=0\}$, with
coefficients of class $C^{m-1,\gamma}$, ellipticity $\eta_{*}$,
obliqueness $\beta_{0}$, and $\norm{v}_{C^{1,\gamma}}\le N_{0}$. For
every $m\ge 2$,
\[
   \norm{v}_{C^{m,\gamma}(B^{+}_{R_{0}/4})}
      \le \mathcal{S}_{m}\bigl(N,\eta_{*},\beta_{0},N_{0},
                                \norm{\mathrm{coeffs}}_{C^{m-1,\gamma}}\bigr).
\]
\end{theorem}

Setting $\eta_{m}:=\norm{\psi}_{C^{m,\gamma}(\{0\le t\le T_{0}/2\})}$,
Theorem~\ref{thm:lieberman} applied iteratively on shrinking half-cylinders
gives
\begin{equation}\label{eq:eta-recursion}
   \eta_{m+1}\le\mathcal{S}_{m+1}\bigl(N,\eta_{*},1,\eta_{m},L_{q}\bigr)
                =:F_{m}(\eta_{m},L_{q},N,R),
\end{equation}
and inductively
\begin{equation}\label{eq:Mm}
   \boxed{\;M_{m}(N,R):=F_{m-1}\bigl(\cdots F_{1}(\eta_{1},L_{q},N,R)
                                            \cdots,L_{q},N,R\bigr),\;}
\end{equation}
with $\eta_{1}=C(N,R)$ from \S\ref{ss:hodo-norm}.

\begin{proposition}[Uniform Schauder bootstrap]\label{prop:bootstrap}
For every $m\ge 1$ there is $M_{m}(N,R)<\infty$, given by~\eqref{eq:Mm},
such that
\begin{equation}\label{eq:bootstrap}
   \norm{g}_{C^{m,\gamma}(\Sph)}\le M_{m}(N,R).
\end{equation}
\end{proposition}

\section{Hölder interpolation to small $C^{2,\gamma_{0}}$}\label{sec:interp}

\begin{lemma}[Hölder interpolation on $\Sph$]\label{lem:interp}
Fix $0\le k<m$, $0<\gamma_{0}\le\gamma\le1$ with $(k,\gamma_{0})\neq(m,\gamma)$,
and set
\begin{equation}\label{eq:theta}
   \theta_{m,k}:=\frac{k+\gamma_{0}}{m+\gamma}\in(0,1).
\end{equation}
There is $C_{\mathrm{GN}}=C_{\mathrm{GN}}(N,m,k,\gamma,\gamma_{0})>0$ such
that for every $f\in C^{m,\gamma}(\Sph)$,
\begin{equation}\label{eq:interp}
   \norm{f}_{C^{k,\gamma_{0}}(\Sph)}
   \le C_{\mathrm{GN}}\,
       \norm{f}_{L^{\infty}(\Sph)}^{1-\theta_{m,k}}\,
       \norm{f}_{C^{m,\gamma}(\Sph)}^{\theta_{m,k}}.
\end{equation}
\end{lemma}

This is the standard GT/BL convexity inequality on a compact manifold;
see \cite[Lemma~6.32]{GT} and \cite[Theorem~6.4.5]{BL}.

We apply Lemma~\ref{lem:interp} with $k=2$, $m=5$:
\[
   \theta:=\theta_{5,2}=\frac{2+\gamma_{0}}{5+\gamma}.
\]
Since $\gamma_{0}\le\gamma\le1$,
$\theta\le(2+\gamma)/(5+\gamma)\le1/2$, hence
\begin{equation}\label{eq:theta-bound}
   \theta\le\tfrac12,\qquad 1-\theta\ge\tfrac12.
\end{equation}

\section{Inversion: closed-form $\alpha_{\sph}(N,R,\gamma_{0})$}\label{sec:closed-form}

\begin{theorem}[Schauder--interpolation entry, quantitative]\label{thm:sph-entry}
Let $0<\gamma_{0}<\gamma\le1$. Choose $m=5$, $\theta=\theta_{5,2}$, and
define
\begin{equation}\label{eq:alpha-sph}
   \boxed{\;
   \alpha_{\sph}(N,R,\gamma_{0})
      :=\left(
              \frac{\delta_{\sph}(N,\gamma_{0})}
                   {C_{\mathrm{GN}}\,C_{H}''(N,R)^{1-\theta}\,
                    M_{5}(N,R)^{\theta}}
            \right)^{\!2N/(1-\theta)}.\;}
\end{equation}
Then every set $\Om$ produced by \texttt{GraphEntry} with
$\alpha(\Om)\le\alpha_{\sph}(N,R,\gamma_{0})$ satisfies
\begin{equation}\label{eq:smallC2}
   \norm{g}_{C^{2,\gamma_{0}}(\Sph)}\le\delta_{\sph}(N,\gamma_{0}),
\end{equation}
the smallness hypothesis of BDV Theorem~3.3.
\end{theorem}

\begin{proof}
By Proposition~\ref{prop:bootstrap} with $m=5$ and
Lemma~\ref{lem:interp} with $k=2$,
\[
   \norm{g}_{C^{2,\gamma_{0}}}
      \le C_{\mathrm{GN}}\,\norm{g}_{L^{\infty}}^{1-\theta}\,
           \norm{g}_{C^{5,\gamma}}^{\theta}
      \le C_{\mathrm{GN}}\,M_{5}(N,R)^{\theta}\,
           \norm{g}_{L^{\infty}}^{1-\theta}.
\]
Insert~\eqref{eq:input-b}:
\[
   \norm{g}_{L^{\infty}}^{1-\theta}
      \le C_{H}''(N,R)^{1-\theta}\,\alpha(\Om)^{(1-\theta)/(2N)}.
\]
$\norm{g}_{C^{2,\gamma_{0}}}\le\delta_{\sph}$ is equivalent to
$\alpha(\Om)^{(1-\theta)/(2N)}\le
\delta_{\sph}/(C_{\mathrm{GN}}C_{H}''^{1-\theta}M_{4}^{\theta})$;
raising to the power $2N/(1-\theta)$ gives $\alpha(\Om)\le\alpha_{\sph}$.
\end{proof}

\begin{remark}[Sharp exponent]
By~\eqref{eq:theta-bound}, $1-\theta\ge1/2$, so the exponent
$2N/(1-\theta)\le 4N$ is uniformly bounded in $\gamma,\gamma_{0}$. The
$L^{\infty}$-power $1-\theta\ge1/2$ recovers the sharp $1/2$ exponent
needed downstream.
\end{remark}

\begin{remark}[Why $m=5$]
Any $m\ge 3$ with $\theta_{m,2}<1$ works for interpolation; $m=5$ makes
$\theta\le1/2$ explicit and uniform for all $0<\gamma_0<\gamma\le1$.
\end{remark}

\begin{remark}[No compactness]
Every step is constructive. $M_{4}(N,R)$ in~\eqref{eq:Mm} is a finite
iteration of Theorem~\ref{thm:lieberman}, whose constants are explicit
in the input data. $C_{\mathrm{GN}}$ is the constant in classical Hölder
interpolation. No sequence, no limit, no contradiction.
\end{remark}

\section{Summary of constants}\label{sec:summary}

\begin{center}
\renewcommand{\arraystretch}{1.20}
\begin{tabular}{@{}lll@{}}
\toprule
Symbol & Origin & Depends on \\
\midrule
$\gamma$ & Alt--Caffarelli (BDV T4.18) & $N,R$ \\
$\gamma_{0}\in(0,\gamma)$ & free parameter & user \\
$L_{u},r_{d}$ & BDV L4.12 & $N,R$ \\
$q_{\pm},L_{q}$ & BDV L4.16 & $N,R$ \\
$M_{1}(N,R)$ & \texttt{GraphEntry}+BDV~4.16 & $N,R$ \\
$C_{H}''(N,R)$ & \texttt{GraphEntry} & $N,R$ \\
$\eta_{*}(N,R)$ & ellipticity of~\eqref{eq:psi-PDE} & $N,R$ \\
$\beta_{0}=1$ & obliqueness~\eqref{eq:BC-psi} & universal \\
$\mathcal{S}_{m}$ & Lieberman~\cite{Li86} Thm.~1 & inputs \\
$F_{m}$ & recursion~\eqref{eq:eta-recursion} & $N,R,L_{q}$ \\
$M_{m}(N,R)$ & iteration~\eqref{eq:Mm} & $N,R$ \\
$\theta_{m,k}=(k+\gamma_{0})/(m+\gamma)$ & GT/BL & $N,m,k,\gamma,\gamma_{0}$ \\
$\theta=\theta_{5,2}\le1/2$ & $m=5,k=2$ & $\gamma,\gamma_{0}$ \\
$C_{\mathrm{GN}}$ & GT/BL interpolation & $N,m,k,\gamma,\gamma_{0}$ \\
$\delta_{\sph}(N,\gamma_{0})$ & BDV T3.3 & $N,\gamma_{0}$ \\
$c_{\sph}(N)$ & BDV T3.3 & $N$ \\
\midrule
$\boxed{\alpha_{\sph}(N,R,\gamma_{0})}$
   & Eq.~\eqref{eq:alpha-sph}
   & $N,R,\gamma_{0}$ \\
\bottomrule
\end{tabular}
\end{center}

\begin{thebibliography}{9}
\bibitem{KN} D.~Kinderlehrer, L.~Nirenberg,
\emph{Regularity in free boundary problems},
Ann.~Scuola Norm.~Sup.~Pisa Cl.~Sci.\ (4) \textbf{4} (1977), 373--391.
\bibitem{Li86} G.~M.~Lieberman,
\emph{The Perron process applied to oblique derivative problems},
Adv.\ in Math.\ \textbf{55} (1985), 161--172; and
\emph{Mixed boundary value problems\ldots},
J.~Math.\ Anal.\ Appl.\ \textbf{113} (1986), 422--440.
\bibitem{GT} D.~Gilbarg, N.~S.~Trudinger,
\emph{Elliptic Partial Differential Equations of Second Order},
2nd ed., Springer, 1983, Lemma~6.32.
\bibitem{BL} J.~Bergh, J.~L\"ofstr\"om,
\emph{Interpolation Spaces. An Introduction},
Springer, 1976, Theorem~6.4.5.
\bibitem{BDV} L.~Brasco, G.~De~Philippis, B.~Velichkov,
\emph{Faber--Krahn inequalities in sharp quantitative form},
Duke Math.\ J.\ \textbf{164} (2015), 1777--1831 (arXiv:1306.0392).
\end{thebibliography}

\end{document}
